% =====================================================================================================
%  Radiology: Artificial Intelligence -- Original Research
%  SHORT, FOCUSED PAPER. Not a cut-down of paper/rsnai/main.tex: that paper asks what a classifier can
%  name; this one asks what survives contact with an operating room.
%
%  Every number in this file is transcribed from the generated evidence briefs, which are themselves
%  written from the result JSONs:
%      docs/COUNTING_EVIDENCE.md          counting arms
%      docs/WRONG_LEVEL_COST_EVIDENCE.md  rates, litigation, fiducial cost
%      docs/WRONG_SITE_BURDEN_EVIDENCE.md annual burden (extrapolations labelled there and here)
%      docs/SWEEP_TIME_EVIDENCE.md        C-arm move times, exposure times, competitor latency
%      docs/SWEEP_PROTOCOL.md             the protocol handed to a radiographer
%  Nothing is typed from memory. A figure absent from those files is stated in the text as not measured.
% =====================================================================================================
\documentclass[11pt,letterpaper]{article}
\usepackage[margin=1in]{geometry}
\usepackage{setspace}\doublespacing
\usepackage[T1]{fontenc}
\usepackage{helvet}
\renewcommand{\familydefault}{\sfdefault}
\usepackage{graphicx,booktabs,array,amsmath,amssymb,xcolor,lineno}
\usepackage[expansion=false]{microtype}
\usepackage[labelfont=bf,font=small]{caption}
\usepackage[numbers,sort&compress]{natbib}
\usepackage[hidelinks]{hyperref}
\linenumbers

\begin{document}

% ------------------------------------------------------------------------------------------------ TITLE
\begin{titlepage}
\begin{center}
{\Large\bfseries Making the Intraoperative Count Exact: Anchored Overlapping Fluoroscopy for Vertebral
Level Confirmation Without a Learned Component\par}
\vspace{1em}
{\normalsize Gregory Schwing, MD, PhD\textsuperscript{1}\par}
\vspace{0.5em}
{\small \textsuperscript{1}Department of Surgery, Detroit Medical Center / Wayne State University School of
Medicine, Detroit, MI\par}
\vspace{1em}
{\small Corresponding author: G. Schwing, [e-mail withheld]\par}
\end{center}
\vspace{1em}
\noindent\textbf{Article type:} Original Research\\
\textbf{Manuscript word count (Introduction--Discussion):} $\approx$2,900\\
\textbf{Tables:} 3 \quad \textbf{Figures:} 0 \quad \textbf{Funding:} none declared.\\
\textbf{Conflicts of interest:} the author has a commercial interest in a level-confirmation product under
development.\\
\textbf{Data sharing:} all imaging is public and de-identified. Every counting arm reported here is written
by one tool into a result file, and the table of arms is generated from those files rather than typed.
\end{titlepage}

% ------------------------------------------------------------------------------------------------ ABSTRACT
\section*{Abstract}

\noindent\textbf{Background:} The Universal Protocol verifies a mark. A vertebral level carries no mark: it
is counted intraoperatively across several limited fluoroscopic fields, and each join between fields can be
one body off. Chart review put the protocol's ceiling at 62\,\% of wrong-site cases prevented, and a
34,544-case registry put wrong-level incidence at 0.12\,\%, rising to 0.42\,\% in the thoracic spine.

\noindent\textbf{Purpose:} To measure whether overlapping exposures linked from an anatomical anchor make
the intraoperative count exact, and to locate the residual risk that remains when it is.

\noindent\textbf{Materials and Methods:} Retrospective study of public labelled CT and radiograph
collections. Frames spanning a 9-inch image-intensifier field (14\,cm at the spine) were cut from each
labelled long image by a fixed protocol---anchor at the edge of the first field, then a move of one field
minus two vertebral bodies---and degraded independently. Consecutive frames were linked by normalised
cross-correlation restricted to within 0.6 of a vertebral body of the step the protocol commanded. The
target is the $k$-th node from the anchor. No learned component takes part in the count and no level name
is formed, so no spine was withheld. Walks were run with the annotation's vertebra boxes and, separately,
with a detector's boxes. The endpoint is the share of walks putting the mark on the body truly $k$ from the
anchor, among walks whose target entered a field before the image ran out.

\noindent\textbf{Results:} With the annotation's boxes, counting down from C2 at three bodies of overlap
was exact on 834 of 834 walks across 76 whole-spine studies at a mean of 2.24 exposures, and on 2,646 of
2,654 reached walks (99.70\,\%) across 204 studies from an institution none of this work trained on.
Counting up from the sacrum was exact on 6,234 of 6,255 reached walks (99.66\,\%) across all 801
abdominopelvic spines and on 7,168 of 7,228 (99.17\,\%) across 522 studies. On 16 real lateral radiographs
reaching the sacrum, every walk was exact (95 of 95) at 1.61 exposures. The residual risk is in the boxes.
On 79 spines held out from a render-trained detector, the annotation's boxes gave 618 of 618 reached walks
exact; the detector's own boxes gave 187 of 622 (30.06\,\%), and the node the count started from was the
true sacrum on 42.44\,\% of walks.

\noindent\textbf{Conclusion:} Counting across overlapping fields is exact and needs no network. What needs
a network is naming, and naming is where the uncertainty lives.

\medskip
\noindent\textbf{Summary Statement:} Overlapping fluoroscopic exposures linked from an anchor made the
intraoperative vertebral count exact on 99.17--100\,\% of reached walks in nine arms of 16 to 801 studies
each, using the annotation's vertebra boxes and with no learned component in the count; with a detector's
boxes on the same spines it fell to 30.06\,\%, placing the whole residual risk in anchor detection.

\clearpage

% ------------------------------------------------------------------------------------------------ INTRO
\section*{Introduction}

The Universal Protocol works by marking. The site is marked while the patient is awake, the mark is
confirmed at the time-out, and the operation proceeds where the mark is. A vertebral level has no mark. It
lies under several centimetres of tissue, it is indistinguishable from its neighbours, and it is identified
after the incision, by counting on fluoroscopy from whatever landmark is visible
\citep{strong2021,mayer2008}. That is not a lapse in discipline; it is a step the protocol does not cover.
Chart review of insured wrong-site cases judged that the Universal Protocol would have prevented eight of
13 under ideal conditions \citep{kwaan2006}, wrong-site events continued at a steady rate through the
protocol era \citep{stahel2010}, and the only systematic review to assess site-verification protocols in
the spine found they failed to prevent about one-third of cases \citep{longo2012}. The remedy the spine
literature recommends---imaging after exposure against a marked fixed structure \citep{devine2010}---is
itself another count.

The failure is rare per procedure and common per career. A prospectively collected registry of 34,544
degenerative spine operations recorded wrong-level surgery in 0.12\,\% (95\,\% CI 0.09--0.16), with the
thoracic spine highest at 0.42\,\% (CI 0.16--1.08), three times the lumbar and eight times the cervical
rate \citep{furuya2026}; anonymous self-report gives one in 3,110 procedures \citep{mody2008}; and about
half of responding neurosurgical spine surgeons in two independent surveys report at least one wrong-level
operation in their career \citep{mody2008,groff2013}. Seventeen percent of self-reported events led to
legal action or settlement \citep{mody2008}; of 68 litigated wrong-disc-level cases, 55 ended in a payment
and surgeons prevailed in 42\,\% of those a jury decided \citep{goodkin2004}; the median indemnity for a
surgical never event reported to the National Practitioner Data Bank is \$75,220 (IQR
\$28,847--\$215,003, 2025 dollars) \citep{bettis2026}; 978 English spinal-surgery negligence claims over
five years were valued at \pounds535.5 million \citep{machin2018}.

Aggregate burden is bounded rather than known, and this paper declines to overstate it. The Joint
Commission received 127 wrong-surgery reports in 2024, nine of them wrong-level and seven of those spinal,
from a voluntary system it states supports no frequency conclusion \citep{jc2024}; NHS England recorded 185
wrong-site never events in one year \citep{nhse2025}. Claims-based modelling puts surgical never events near
4,000 a year in the United States \citep{mehtsun2013}, and applying the two per-procedure rates above to a
United States volume of 1.3--1.6 million spine operations gives roughly 400 to 1,900 wrong-level operations
and, with the verified litigation and payment fractions and a mean paid claim near \$133,000--\$136,000
\citep{mehtsun2013,tan2023}, on the order of \$8--37 million a year in indemnity alone. Each step of that is
an extrapolation, the spine-surgery denominator is unverified, and hospital and disability costs have no
published per-event figure. The often-quoted \$1.3 billion is a 21-year cumulative for all surgical never
events \citep{mehtsun2013}, not an annual wrong-site figure, and no defensible worldwide figure exists.

What is spent avoiding the error is easier to price. One quaternary centre placed 179 preoperative
percutaneous gold fiducials over 38 months, 176 of them thoracic, at a mean of 51.9 minutes and 183\,mGy
each, purely to mark a level for the surgeon \citep{siminski2023}. At a hospital operating-room cost of
\$34--37 per minute before physician fees \citep{childers2018}, that is roughly \$1,750--\$1,950 of staffed
procedural time per patient and about 155 hours across the series---a proxy, since the marker is placed in a
fluoroscopy or CT suite rather than a theatre. The fiducial buys real time: thoracic localisation falls from
about 15 minutes to about three \citep{mongardi2021,upadhyaya2012}, at the cost of a separate 14.3-minute
CT intervention \citep{keil2024}.

Automated alternatives have been built almost entirely on registering a preoperative CT to the intraoperative
radiograph \citep{levelcheck2012,otake2015}. Registration does not name a vertebra. It solves a
six-degree-of-freedom pose and transfers a name a human read off the preoperative volume, which means it is
unavailable when there is no CT, and a miscounted CT propagates onto the radiograph faithfully.

This paper starts from what surgeons actually do. They do not classify a vertebra by its shape; they count,
from C2 down or the sacrum up, and the count is unreliable only because it is carried across several
limited fields whose joins are ambiguous in a nearly periodic column. The intervention that follows is not
a classifier that asserts an identity but a mechanism that makes the join exact. We measure how exact,
where it breaks, and what it costs to run.

% ------------------------------------------------------------------------------------------------ METHODS
\section*{Materials and Methods}

This retrospective study used public, de-identified imaging only; institutional review determined it not to
be human-subjects research.

\subsection*{Imaging}
Four sources were used. CTSpinoPelvic1K supplied 801 abdominopelvic CT studies with a per-level vertebral
label \citep{ctsp1k}; VerSe supplied 76 whole-spine studies \citep{verse2021}; TotalSegmentator's public
release of routine clinical CT supplied 204 studies holding C2 and 522 reaching the sacrum, from an
institution on which nothing in this work has been trained; and real lateral radiographs supplied 16 spines
reaching the sacrum and 24 holding C2. Serial sweeps of one patient do not exist in any public collection,
so each sweep was cut from a single labelled long image. Every frame was degraded independently---gain,
offset, gamma, a small rotation and scale, and a fitted film-noise model---so no two frames share pixels
exactly.

\subsection*{The sweep protocol}
The protocol is the one a radiographer already performs by hand, made explicit \citep{schwing2026ctx}. The
anchor is placed at the edge of the first field: the sacrum at the bottom for lumbar and thoracolumbar work,
C2 at the top for cervical and upper thoracic work. The C-arm then moves by one field minus two vertebral
bodies, and the exposure is repeated to the level of interest. The last move is a full move or none. Frame
span is physical rather than a count of vertebrae: a 9-inch image intensifier spans about 14\,cm at the
spine once the fixed 100\,cm source-to-detector distance of a mobile C-arm is accounted for
\citep{naik2022,ge_k170752}, which holds about four lumbar bodies.

\subsection*{Linking, and the count}
Consecutive frames are linked by normalised cross-correlation on the overlap band alone, never by a global
homography, and the search is restricted to within 0.6 of a vertebral body of the step the protocol
commanded. That restriction is the mechanism of the paper. An unrestricted correlation on a nearly periodic
column locks one body off, which is the documented failure mode of automatic vertebral localisation and
accounted for all 47 failures in the original three-dimensional to two-dimensional registration study
\citep{otake2012fail}. Restricting the search to the commanded displacement removes the aliasing, because
the protocol states what the answer should be to within half a body.

The linked frames form one chain of vertebral nodes. A walk starts at the anchor node and the target is the
$k$-th node from it. Nothing is fitted anywhere in this procedure: no encoder is loaded, no level name is
formed, no alphabet is chosen, and variant anatomy therefore never enters. Because nothing is fitted, no
spine is withheld on a training-distribution argument, and every study carrying the anchor was used.

\subsection*{Box sources}
The chain's nodes are vertebra boxes. Headline arms use the annotation's boxes, which isolates the linker:
the only thing that can fail is the join. A second set of arms substitutes a detector's boxes and is
restricted to spines held out from that detector. In the arms reported here the anchor node is the terminal
node of the chain at the anchor end; identifying it by class instead is a separate arm that has not
returned.

\subsection*{Endpoints}
\emph{Reached} is the number of walks whose target entered a field before the image ran out; a walk that
does not reach is a limitation of field, not a counting error, and both denominators are reported.
\emph{Exact of reached} is the share of reached walks placing the mark on the body whose position in the
labels is truly $k$ from the anchor. \emph{Exposures} is the mean number of frames a walk consumed.
\emph{Anchor held} is a diagnostic---how often the node the count began from really was the anchor---and is
never used as a gate. \emph{Refused} counts walks the system declined. Wrong walks are reported with the
number of distinct spines they fall in, because a single bad study produces many bad walks.

\subsection*{Naming, and timing}
For the division-of-labour comparison, the same chains were named by a ResNet-18 encoder over fixed-scale
vertebra crops with a bidirectional recurrent head over the chain \citep{he2016,schwing2026mixed}. Per-frame
encoder time, per-walk linking time and per-walk decode time were measured in the same runs. No hypothesis
test was applied; proportions are reported with their numerators and denominators.

% ------------------------------------------------------------------------------------------------ RESULTS
\section*{Results}

\subsection*{The count is exact, and it needs no network}
Table~\ref{tab:count} gives every counting arm. With the annotation's boxes, counting down from C2 at three
bodies of overlap was exact on 834 of 834 walks over 76 whole-spine studies, at a mean of 2.24 exposures;
at two bodies, 832 of 834, the two failures in one study. On 204 TotalSegmentator studies holding C2 the
same protocol was exact on 2,646 of 2,654 reached walks (99.70\,\%) at two bodies and 2,640 of 2,650
(99.62\,\%) at three, the misses confined to two spines in each arm. Counting up from the sacrum was exact
on 6,234 of 6,255 reached walks (99.66\,\%) over all 801 abdominopelvic spines at two bodies of overlap and
6,214 of 6,229 (99.76\,\%) at three, with errors in six and five spines; on 522 TotalSegmentator spines,
7,168 of 7,228 (99.17\,\%), the errors in six spines. Every arm consumed between 1.15 and 2.96 exposures.

The errors are not scattered. They sit far from the anchor. Counting up from the sacrum on the abdominopelvic
spines was exact on 99.73\,\% of lumbar targets and 95.45\,\% of thoracic ones at two bodies of overlap, and
on the TotalSegmentator sacrum arm the cervical targets---18 or 19 bodies from the anchor---fell to
83.84\,\% against 99.33\,\% lumbar. Counting down from C2 on the whole-spine studies shows the same
gradient from the other end: 99.46\,\% thoracic at two bodies of overlap and 100\,\% at three. A count degrades with the number of joins
it crosses, exactly as a surgeon's does; the remedy is to anchor at the near end, which the protocol permits
because both ends of the column carry an unmistakable landmark.

What did not reach is reported separately. Ninety-seven of 6,352 walks on the abdominopelvic spines and 168
of 7,396 on TotalSegmentator ended before the target appeared, almost all of them the topmost labelled body
or a cervical level far from the sacrum. Those are field limitations and are not counted as errors.

\subsection*{On real radiographs}
On 16 real lateral radiographs reaching the sacrum, every walk was exact with the annotation's boxes: 95 of
95 reached, at 1.61 exposures. On 24 real radiographs holding C2, 110 of 110, at 1.15 exposures. Sixteen and
24 spines are thin, and the number is stated beside the percentage for that reason. These two arms establish
that the linker survives real radiographic pixels; they do not establish a rate.

\subsection*{The residual risk is the anchor, not the count}
Table~\ref{tab:boxes} is the result that governs deployment. On 79 CTSpinoPelvic1K spines held out from a
render-trained detector, the annotation's boxes gave 618 of 618 reached walks exact. The detector's own
boxes on the same spines gave 187 of 622 (30.06\,\%), with errors in all 79 spines; at three bodies of
overlap, 178 of 621 (28.66\,\%) in 75 spines. The mechanism is visible in the diagnostic: the node the count
started from was the true sacrum on 42.44\,\% of walks at two bodies and 39.77\,\% at three. A missed or
spurious box shifts the ordinal, and shifting the ordinal renames every body above it. That is the surgeon's
own error, reintroduced by a machine, and it is why the count alone is not a product.

The same 16 real films run with the deployed film detector's boxes gave 95 of 95 exact walks at 1.47
exposures while the anchor diagnostic read 93.68\,\%: some walks were scored as beginning from a node that
was not the true sacrum and nonetheless ended on the correct body. Those two numbers are not reconciled
here. Until they are, that arm is not evidence that a detector's boxes are safe; the 79-spine CT arm is the
arm to believe.

Three arms address this directly and none has returned. Identifying the anchor by class rather than by
chain position, on the CTSpinoPelvic1K detector boxes at two and at three bodies of overlap, and on the real
films, is \textbf{[PENDING: anchor-by-class counting arms]}. A detector holding an explicit C2 class is
\textbf{[PENDING: C2-class anchor detector]}. Confirming a level by re-imaging with a radiopaque instrument
in the field and placing that instrument on the chain is \textbf{[PENDING: needle confirmation]}; the
evidence brief carries the heading and no result.

One further gap is reported rather than explained: the refusal count is zero in every arm in
Table~\ref{tab:count} and Table~\ref{tab:boxes}. The system's declared refusal behaviour---declining an
uncertain link, or a later exposure it cannot place---never fired in these runs. Refusal is therefore
specified and unmeasured, and a system whose safety argument rests on refusing when it cannot see the anchor
has not yet demonstrated that it refuses.

\subsection*{Naming is the part that is uncertain}
The same chains named by the encoder and its chain decoder were correct on 96.25\,\% of 950 walks on
CTSpinoPelvic1K and 94.69\,\% of 7,405 walks on TotalSegmentator, the latter corpus never trained on
anywhere in this work. By region on TotalSegmentator: 96.21\,\% lumbar, 91.58\,\% thoracic, 80.35\,\%
cervical. The evidence brief does not state how many distinct spines the 950 CTSpinoPelvic1K walks span,
and that is not recorded here.

Those figures sit four to five points below the counting arms measured on the same protocol, which is the
paper's central comparison: the learned step is the worse step.

\subsection*{What it costs to run}
Table~\ref{tab:time} gives the budget. The encoder takes 17--18\,ms per frame, linking a whole walk takes
0.102--0.29\,s, and decoding takes 0.2\,ms. Against that, repositioning a mobile C-arm takes 36--49\,s per
move in the only published time-motion measurements \citep{matthews2007,desilva2018}, and beam-on time per
view is 0.5--4.7\,s, about 11\,\% of positioning time \citep{desilva2018}. A two-exposure walk needs one
move beyond the first placement. The compute is not the constraint and no speed claim is made: the
comparison a reviewer should apply is against 4.7\,s to 72\,s per image for the clinically translated
registration system \citep{levelcheck2012,manbachi2018}, against which surgeons said they would wait 20\,s
to one minute \citep{desilva2016}.

% ------------------------------------------------------------------------------------------------ DISCUSSION
\section*{Discussion}

The result that matters is a division of labour. Counting across overlapping fields is exact and contains no
learned component: 99.17--100\,\% of reached walks in nine arms of 16 to 801 studies each, at 1.15 to 2.96
exposures, with the annotation's boxes. Naming the same chains with a trained network was 94.7--96.3\,\%. The two numbers come
from the same images, the same protocol and the same run, and the learned step is the weaker one. Converting
an ordinal into a name requires an alphabet, and transitional anatomy changes the alphabet; a network's
proper job here is not to recognise each vertebra from its morphology but to decide which alphabet the
patient has. When a surgical plan is stated relative to the anchor---the third mobile segment above the
sacrum rather than L4---counting alone suffices, and nothing needs to be learned at all.

\noindent\textbf{What has to be true in the room.} Three things. The anchor must be in the first field,
which the radiographer controls. The vertebrae must be found in every frame, which a detector controls and
which is where all of the residual risk now sits: substituting a detector's boxes for the annotation's
dropped exactness from 100\,\% to 30.06\,\% on the same 79 spines, and the anchor node was correct on fewer
than half the walks. And the patient must not move between exposures, which is not measured here at all.

\noindent\textbf{What happens when it fails, and what it refuses.} A wrong join or a wrong anchor shifts the
count by one body and renames the column above it confidently. The mitigation designed for this is
refusal---identify the anchor by class and decline when no anchor is found---and, as reported above, refusal
fired zero times in every arm, so it is a design and not a result. The second mitigation is that the failure
is auditable by eye: the chain is a strip of overlapping radiographs, and a surgeon can see whether the
bodies shared between two frames are the same bodies. That is a property registration does not have, because
a registration failure looks like a correctly placed overlay.

\noindent\textbf{What it does not do.} It forms no level name, so it cannot tell a surgeon that a body is L5
rather than L6; that question belongs to the naming model and is open. It does not reach targets the image
does not cover, and 97 and 168 walks in two corpora ended for that reason. It was not measured on a moving
patient, on parallax between exposures taken at different gantry angles, or in any operating room; every
sweep in this paper was cut from a single long image. There is no reader study and no prospective arm.

\noindent\textbf{Why this is deployable where navigation is not.} The mechanism needs an anchor, two to
three exposures on a C-arm the hospital already owns, and an embedded computer: 17--18\,ms per frame is not a
workstation-class load. It needs no preoperative CT, no optical tracker and no line of sight, and it asks for
one model rather than two, because the count requires none. This is the setting in which wrong-level surgery
is least protected---community hospitals and ambulatory surgery centres that do not buy navigation
suites---but the capital cost of those suites was not measured in this study and no saving is claimed from
avoiding one. Against the alternative that is priced, a preoperative gold fiducial costs 51.9 minutes and
183\,mGy per patient \citep{siminski2023}, and a two-exposure sweep costs one C-arm move.

\noindent\textbf{Where it belongs in the Universal Protocol.} The protocol verifies a mark; this supplies the
spine's missing one. Verification enters the surgical field with the patient, the sweep places a digital
mark on the target, and the time-out compares that mark's count against the surgeon's own. It is assistance
and not replacement: two independent counts that must agree, with the disagreement displayed. A quality
programme can audit it, because the evidence is a saved strip of images rather than a number a black box
emitted.

\noindent\textbf{Limitations.} Every headline number uses the annotation's vertebra boxes, which measures the
linker and not a deployable system; the detector-box arms are the honest preview and they are poor. The
anchor-by-class arms that are meant to fix them have not returned. Sweeps are simulated from single long
images and reproduce independent exposure, noise, rotation and scale between frames but not parallax and not
patient motion, both of which act on linking. The correlation search is restricted to within 0.6 of a body
of the commanded step, and in simulation the step taken is the step commanded, by construction; a
radiographer's move is not, and the width of that discrepancy has not been measured. The real-radiograph arms rest on 16 and 24 spines. Refusal
never fired. Two corpora---CTSpinoPelvic1K and VerSe---overlap with the training data of the naming model
reported for comparison, although not of the count, which is fitted on nothing. TotalSegmentator carries no
variant lumbar or thoracic class, so nothing here bears on transitional anatomy. The cost figures in the
Introduction are extrapolations from verified rates onto an unverified surgical volume, and are stated as
ranges for that reason.

% ------------------------------------------------------------------------------------------------ TABLES
\clearpage
\section*{Tables}

\begin{table}[htbp]\centering\small
\caption{Counting arms with the annotation's vertebra boxes. No learned component takes part, so no spine is
withheld. \emph{Reached} excludes walks whose target never entered a field. \emph{Wrong} gives walks with
distinct spines in parentheses. The anchor diagnostic was 1.0000 in every row except the 522-spine sacrum
arm (0.9935). Refusals were zero in every row. Transcribed from \texttt{docs/COUNTING\_EVIDENCE.md}.}
\label{tab:count}
\begin{tabular}{llrrrrrr}
\toprule
Anchor & Corpus & Spines & Walks & Reached & Exact & Exp. & Wrong \\
\midrule
C2, 2 bodies      & VerSe               & 76  & 834   & 834   & 99.76\,\%  & 1.99 & 2 (1) \\
C2, 3 bodies      & VerSe               & 76  & 834   & 834   & 100\,\%    & 2.24 & 0 (0) \\
C2, 2 bodies      & TotalSegmentator    & 204 & 2,656 & 2,654 & 99.70\,\%  & 2.05 & 8 (2) \\
C2, 3 bodies      & TotalSegmentator    & 204 & 2,656 & 2,650 & 99.62\,\%  & 2.27 & 10 (2) \\
C2, 2 bodies      & Real radiographs    & 24  & 110   & 110   & 100\,\%    & 1.15 & 0 (0) \\
\midrule
Sacrum, 2 bodies  & CTSpinoPelvic1K     & 801 & 6,352 & 6,255 & 99.66\,\%  & 1.98 & 21 (6) \\
Sacrum, 3 bodies  & CTSpinoPelvic1K     & 801 & 6,352 & 6,229 & 99.76\,\%  & 2.42 & 15 (5) \\
Sacrum, 2 bodies  & TotalSegmentator    & 522 & 7,396 & 7,228 & 99.17\,\%  & 2.96 & 60 (6) \\
Sacrum, 2 bodies  & Real radiographs    & 16  & 95    & 95    & 100\,\%    & 1.61 & 0 (0) \\
\bottomrule
\end{tabular}
\end{table}

\begin{table}[htbp]\centering\small
\caption{Where the residual risk is. The first three rows are the same 79 CTSpinoPelvic1K spines, held out
from the detector, counted from the sacrum with three different box sources. \emph{Anchor} is the share of
walks whose starting node really was the sacrum. The film row's exactness and anchor diagnostic are not
reconciled in this study. Transcribed from \texttt{docs/COUNTING\_EVIDENCE.md}.}
\label{tab:boxes}
\begin{tabular}{llrrrrr}
\toprule
Boxes & Overlap & Spines & Reached & Exact & Anchor & Wrong \\
\midrule
Annotation        & 2 bodies & 79 & 618 & 100\,\%   & 100\,\%   & 0 (0) \\
Render detector   & 2 bodies & 79 & 622 & 30.06\,\% & 42.44\,\% & 435 (79) \\
Render detector   & 3 bodies & 79 & 621 & 28.66\,\% & 39.77\,\% & 443 (75) \\
\midrule
Film detector     & 2 bodies & 16 & 95  & 100\,\%   & 93.68\,\% & 0 (0) \\
\bottomrule
\end{tabular}
\end{table}

\begin{table}[htbp]\centering\small
\caption{The deployment budget. Rows without a citation were measured in this study's own runs. The compute
is three orders of magnitude below one C-arm move, so the argument is not speed; it is that the count needs
no model and the hardware is a C-arm the hospital already owns.}
\label{tab:time}
\begin{tabular}{lll}
\toprule
Quantity & Value & Source \\
\midrule
Encoder, per frame (naming only)     & 17--18\,ms          & this study \\
Frame linking, per walk              & 0.102--0.29\,s      & this study \\
Chain decode, per walk               & 0.2\,ms             & this study \\
Exposures per walk                   & 1.15--2.96          & this study \\
\midrule
C-arm reposition, per move           & 36--49\,s           & \citep{matthews2007,desilva2018} \\
Beam-on, per view                    & 0.5--4.7\,s         & \citep{desilva2018} \\
Registration labelling, per image    & 4.7--72\,s          & \citep{levelcheck2012,manbachi2018} \\
Surgeon-stated acceptable wait       & 20\,s--1\,min       & \citep{desilva2016} \\
\midrule
Thoracic localisation, fluoroscopy   & $\approx$15\,min    & \citep{mongardi2021,upadhyaya2012} \\
Thoracic localisation, with fiducial & $\approx$3\,min     & \citep{mongardi2021,upadhyaya2012} \\
Fiducial placement itself            & 51.9\,min, 183\,mGy & \citep{siminski2023} \\
Fiducial CT-suite intervention       & 14.3\,min           & \citep{keil2024} \\
\bottomrule
\end{tabular}
\end{table}

\clearpage
\bibliographystyle{unsrtnat}
\bibliography{../../rsnai/refs}

\end{document}
